% ============================================================
\subsubsection{Pruebas de caja negra}
% ============================================================

Los endpoints de la API se validaron con pruebas de caja negra ejecutadas en
Thunder Client. Cada caso define la petición, el código HTTP esperado y, cuando
aplica, un texto que debe aparecer en la respuesta. La colección completa se
encuentra en el archivo \texttt{NutriApp\_ThunderClient.json}
\pendiente{referencia al anexo}. La Tabla~\ref{tab:pcn} resume los resultados.

\pendiente{ejecuta la colección y llena la columna «Estado» con lo que
realmente obtuviste; anota los casos fallidos y qué se corrigió}

{\footnotesize
\begin{longtable}{|p{1.2cm}|p{2.9cm}|p{5.6cm}|p{2.6cm}|p{1.6cm}|}
\caption{Pruebas de caja negra de la API REST}\label{tab:pcn}\\
\hline
\textbf{ID} & \textbf{Endpoint} & \textbf{Descripción} & \textbf{Esperado} & \textbf{Estado} \\ \hline
\endfirsthead
\hline
\textbf{ID} & \textbf{Endpoint} & \textbf{Descripción} & \textbf{Esperado} & \textbf{Estado} \\ \hline
\endhead
TC-01 & GET /health & Verificar que el servidor está disponible & 200, «ok» & \pendiente{} \\ \hline
TC-02 & POST /api/auth/login & Rechazar inicio de sesión sin cuerpo & 400, «requeridos» & \pendiente{} \\ \hline
TC-03 & POST /api/auth/login & Rechazar contraseña incorrecta & 401, «incorrectos» & \pendiente{} \\ \hline
TC-04 & GET /api/auth/me & Impedir acceso a ruta protegida sin token & 401, «Token requerido» & \pendiente{} \\ \hline
TC-05 & GET /api/auth/me & Rechazar token malformado & 401, «inválido» & \pendiente{} \\ \hline
TC-06 & GET /api/auth/me & Token válido de usuario inexistente en BD & 200 & \pendiente{} \\ \hline
TC-07 & GET /api/patients & Listar pacientes registrados & 200, arreglo & \pendiente{} \\ \hline
TC-08 & POST /api/patients & Rechazar paciente con campos faltantes & 400, «campos requeridos» & \pendiente{} \\ \hline
TC-09 & GET /api/patients/:id & Paciente con ID inexistente & 404, «no encontrado» & \pendiente{} \\ \hline
TC-10 & GET /api/foods & Listar alimentos del catálogo SMAE & 200, arreglo & \pendiente{} \\ \hline
TC-11 & GET /api/foods/smae-groups & Obtener grupos SMAE disponibles & 200 & \pendiente{} \\ \hline
TC-12 & GET /api/foods/:id & Alimento con ID inexistente & 404 & \pendiente{} \\ \hline
TC-13 & POST /api/ai/chat & Rechazar mensaje vacío & 400, «requerido» & \pendiente{} \\ \hline
TC-14 & POST /api/ai/chat & Rechazar petición sin campo \texttt{message} & 400 & \pendiente{} \\ \hline
TC-15 & POST /api/ai/recommendation & Rechazar recomendación con datos del paciente faltantes & 400 & \pendiente{} \\ \hline
TC-16 & POST /api/ai/analyze-food & Rechazar análisis sin campos requeridos & 400 & \pendiente{} \\ \hline
TC-17 & POST /api/ai/meal-plan & Rechazar plan sin calorías objetivo & 400 & \pendiente{} \\ \hline
TC-18 & GET /api/ai/chat-history/global & Obtener historial del chat global & 200, arreglo & \pendiente{} \\ \hline
TC-19 & POST /api/ai/chat & Solicitar cálculo de GET nutricional & 200, «kcal» & \pendiente{} \\ \hline
TC-20 & POST /api/ai/chat & Rechazar solicitud fuera de dominio (código Python) & 200, «Solo puedo asistirte» & \pendiente{} \\ \hline
TC-21 & POST /api/ai/recommendation & Recomendación completa con equivalentes SMAE & 200, «SMAE» & \pendiente{} \\ \hline
TC-22 & POST /api/ai/analyze-food & Análisis de un alimento SMAE (tortilla de maíz) & 200, «kcal» & \pendiente{} \\ \hline
\end{longtable}
}

\figura{im10_postman_rag.png}{Ejemplo de ejecución de una prueba sobre el asistente de IA}{fig:prueba_rag}

\paragraph{Verificación de requisitos no funcionales}

\pendiente{opcional pero recomendable: mide y reporta RNF01 (tiempo de respuesta
del asistente, p.\,ej. promedio de 10 consultas) y RNF03 (compara IMC, ICC y GET
del sistema contra cálculo manual en 3--5 pacientes de prueba)}

\begin{table}[H]
\centering\small
\caption{Verificación de requisitos no funcionales}\label{tab:rnf_verif}
\begin{tabularx}{\linewidth}{|l|X|X|}
\hline
\textbf{RNF} & \textbf{Criterio} & \textbf{Resultado} \\ \hline
RNF01 & Respuesta del asistente $\leq$ 10\,s & \pendiente{tiempo medido} \\ \hline
RNF02 & Contraseñas con bcrypt y sesión con JWT & Verificado en TC-02 a TC-05 \\ \hline
RNF03 & Error $\leq$ 1\,\% en IMC, ICC y GET & \pendiente{error medido} \\ \hline
RNF04 & Funciona en Chrome, Firefox y Edge & \pendiente{} \\ \hline
RNF05 & Ejecución local con Groq como único servicio externo & \pendiente{} \\ \hline
RNF07 & Repositorio con ramas \texttt{main} y \texttt{develop} & \pendiente{} \\ \hline
\end{tabularx}
\end{table}
